% !PRINTED 776
\weberpage{776}
%<<CONT
mit $\Phi(x)$ gemein hat, so muss $\Phi(x)$ durch $f_n(\alpha^{-1}x)$ theilbar sein, und wir können demnach für ein unbestimmtes $x$

\weq{18}{\Phi(x) = f_n(\alpha^{-1}x)\ \Psi(x,\alpha)}

setzen, worin $\Psi(x,\alpha)$ eine ganze rationale Function von $x$ ist.

Sind $\alpha, \alpha_1, \alpha_2, \ldots$ die sämmtlichen primitiven $m^{\text{ten}}$ Einheitswurzeln, die, wie wir angenommen haben, Wurzeln einer \gesperrt{irreduciblen} rationalen Gleichung

\[
F_m(x) = 0
\]

sind, so kann $\alpha$ nach dieser Voraussetzung in der Identität (18) durch jedes andere $\alpha_1, \alpha_2, \ldots$ ersetzt werden, und es folgt, dass $\Phi(x)$ durch jeden der Factoren

\[
f_n(\alpha^{-1}x),\ f_n(\alpha_1^{-1}x),\ f_n(\alpha_2^{-1}x) \ldots
\]

theilbar ist. Zwei dieser Factoren können aber keinen gemeinschaftlichen Theiler haben, denn wenn etwa

\[
f_n(\alpha^{-1}x) = 0,\ f_n(\alpha_1^{-1}x) = 0
\]

eine gemeinsame Wurzel $x$ hätten, so müssten zwei primitive $n^{\text{te}}$ Einheitswurzeln $r, r_1$ existiren, so dass

\[
x = \alpha r = \alpha_1 r_1
\]

wäre; es müsste also $\alpha^n = \alpha_1^n$, und folglich, da $n$ relativ prim zu $m$ ist, $\alpha = \alpha_1$ sein.

Demnach ist $\Phi(x)$ theilbar durch das Product

\[
F_{mn}(x) = f_n(\alpha^{-1}x)\ f_n(\alpha_1^{-1}x)\ f_n(\alpha_2^{-1}x) \ldots,
\]

dessen Wurzeln die sämmtlichen $\beta$ sind, und diese Function ist daher auch unzerlegbar, wie bewiesen werden sollte.

\begin{center}
\rule{4cm}{0.4pt}
\end{center}
